\documentclass[11pt]{beamer}
\usepackage{tikz}
\setbeamertemplate{background canvas}{\begin{tikzpicture}\node[opacity=.1]{\includegraphics
		[width=\paperwidth]{Figuras/Capa/CirgMOdelo.png}};\end{tikzpicture}}
\usetheme{Warsaw}
\usepackage[utf8]{inputenc}
\usepackage[brazil]{babel}  % idioma
\usepackage{amsmath,amsfonts,amssymb,textcomp}
\usepackage{graphicx}
\usepackage{subfigure}
\usepackage{mathtools}
\usepackage[lined]{algorithm2e}


%% Colocando numeração nos Slides


\author{ Othon Luiz Teixeira de Oliveira }
\title[PPGES]{Defesa de Mestrado}
%\setbeamercovered{transparent} 
%\setbeamertemplate{navigation symbols}{} 
%\logo{} 
\institute{Escola Politécnica de Pernambuco -- Poli ---  UPE} 
\date{29 - Maio - 2017} 
%\subject{} 
\begin{document}


\begin{frame}
	\titlepage
\end{frame}


% Capa - requer o TikZ
\newcommand{\capa}{
    \begin{tikzpicture}[remember picture,overlay]
        \node at (current page.south west)
            {\begin{tikzpicture}[remember picture, overlay]
                \fill[shading=radial,top color=orange,bottom color=orange,middle color=yellow] (0,0) rectangle (\paperwidth,\paperheight);
            \end{tikzpicture}
          };
    \end{tikzpicture}
}

%+++++++++++++++++++++++++++++++++++++++++++++++
\begin{frame}{PPGES}
	\begin{figure}[h]
		%\left
		\includegraphics[width=16mm, height=13mm]{Figuras/Capa/upelogo.eps}
		\qquad \quad \quad \quad \quad
		\includegraphics[width=17mm, height=15mm]{Figuras/Capa/logoPoli.jpg}
		\qquad \quad \quad \quad \quad \quad \quad 	\vspace{0.5in}
		\includegraphics[width=20mm, height=15mm]{Figuras/Capa/icblogo.png}
		\\
		\includegraphics[width=55mm, height=35mm]{Figuras/Capa/logo_ppges3.png}
		
	\end{figure}
\end{frame}

\begin{frame}
	\frametitle{ \Large Programa de Pós-Graduação em Engenharia de Sistemas}
	\begin{center}
		\Large \textbf{DISSERTAÇÃO DE MESTRADO} 
	\end{center}
	\pause
	\begin{block}
		\Large \textbf{Modelo Preditivo para Sugestão de Roteamento de cargas considerando dados históricos, sócio-ambientais, e de redes sociais}
	\end{block}
	
	\pause
	\vspace{0.3in}
	%\begin{block}
	\textbf{Mestrando:} Othon Luiz Teixeira de Oliveira \\
	\textbf{Orientador:} Prof. Dr. Fernando Buarque de Lima Neto
	%\end{block}
	
\end{frame}

%%% Colocando o sumário em cada início de Sessão
\AtBeginSection[]{ 
	\begin{frame} 
		\frametitle{Sumário} \tableofcontents[currentsection, hideothersubsections] 
	\end{frame} 
} 
%%% Colocando o sumário em cada início de Subsessão
\AtBeginSubsection[ ]{
	\begin{frame}{Organização}
		\tableofcontents[currentsection,currentsubsection]
	\end{frame}
}


%% Adiciona numeração a partir daqui, porém não parte do zero
\addtocounter{framenumber}{-1}
\setbeamertemplate
 {footline}{\quad\hfill\insertframenumber/\inserttotalframenumber\strut\quad} 

%+++++++++++++++++++++++++++++++++++++++++++++++
\section{ Introdução}
\subsection{Introdução}


\begin{frame}\frametitle{Resumo}	
	\begin{block}{O cenário}
		\begin{itemize}
			\item O transporte de cargas no Brasil é feito principalmente pelas rodovias federais (BRs). 
			\pause
			\item Essas rodovias estão constantemente congestionadas nos perímetros urbanos.
			\pause
			\item Comunidades bloqueiam as rodovias para reivindicar, dos entes públicos, todo tipo de necessidades
			\pause
			\item Em alguns trechos o traçado das rodovias está próximo a morros e florestas.
		\end{itemize}
	\end{block}
\end{frame}



%--------------------------------------------------------
%\section{ Objetivos}
\subsection{ Objetivos}

\begin{frame}\frametitle{Proposição de um modelo preditivo}
	\pause
	\begin{block}{ Objetivo Geral}
		Esta pesquisa teve como principal objetivo desenvolver um modelo de plataforma autoadaptável 
		\pause que contemple predição do comportamento das rodovias federais pernambucanas,
		\pause antecipando alguns eventos que nela possam ocorrer, apontando onde ocorrerão.
		
	\end{block}
\end{frame}

\begin{frame}\frametitle{ Objetivos Específicos}
	\pause
	\begin{block}{ Quatro objetivos}
		\begin{itemize}
			\item Caracterizar a problemática de cada rodovia;
			\pause
			\item Desenvolver um modelo preditivo dos fenômenos que envolvem as rodovias;
			\pause
			\item Desenvolver um ambiente de simulações interativas da estrutura viária em sua dinâmica
			\pause			
			\item Propor soluções para melhorar a experiência dos usuários que utilizam as rodovias pernambucanas
		\end{itemize}
		
	\end{block}
\end{frame}

\section{Background}
\subsection{Dados a minerar}

\begin{frame}\frametitle{Dados originais da PRF}
	%	\begin{block}{PRF & Twitter}
	\begin{itemize}
		\item Base de dados da Polícia Rodoviária Federal: acidente e interdições, entre 2007 a 2015.
		\pause
		\item Rede Social Twitter com 3200 tweets (limite permitido), de 2017 a 2014
	\end{itemize}	
	%	\end{block}
\end{frame}

%++++++++++++++++++++++++++++++++++++++++++++++++++++++++++++++++++++++++++++++++++++++++++++++

\begin{frame}\frametitle{ Planilha com os dados originais}
	\transboxin[duration=1, direction=25]	
	\begin{block}{ Planilha PRF}
		\begin{figure}[!ht]
			\centering %\centering % para centralizarmos a figura
			\caption{Planilha para Preprocessamento}
			\includegraphics[width=110mm, height=40mm]{Figuras/BigData/PlanilhaPRF.png}
		\end{figure}
	\end{block}
\end{frame}


\begin{frame}
	\frametitle{ Entendendo a base de dados da PRF}
	\begin{exampleblock}{ Atributos e Instâncias -- iniciais}
		\begin{itemize}
			\item Dados de 2007 à 2015 -- BRs de Pernambuco
			\pause
			\item 85.209 Instâncias -- 27 Atributos
			\pause
			\item Dentre eles: Km, Latitude, Longitude, Condições da Pista, Causa do Acidente, Município, Data, Hora, Tipo de Veículo, Quantidade de Mortos, etc.
			\pause
			\item 40\% missing data -- 70\% do tempo total para tratar
		\end{itemize}
	\end{exampleblock}
\end{frame}

\begin{frame}\frametitle{ O Tweeter}
	\begin{block}{ Acesso aos dados do Tweeter}
		\begin{figure}[!ht]
			\centering %\centering % para centralizarmos a figura
			\caption{Registro de App para para acessar e baixar dados}
			\includegraphics[width=80mm, height=35mm]{Figuras/BigData/appTweeter.png}
		\end{figure}
	\end{block}
\end{frame}

\begin{frame}\frametitle{ Dados originais}
	\transblindsvertical[duration=2, direction=25]
	\begin{block}{ Dados -- Planilha Tweeter}
		\begin{figure}[!ht]
			\centering %\centering % para centralizarmos a figura
			\caption{Planilha para Preprocessamento}
			\includegraphics[width=110mm, height=40mm]{Figuras/BigData/tweetPRF.png}
		\end{figure}
	\end{block}
\end{frame}

\begin{frame}\frametitle{ Entendendo a base de dados do Tweeter}
	\begin{exampleblock}{Atributos e Instâncias -- iniciais}
		\begin{itemize}
			\item Dados de 2017 -- 2014 --- canal @PRF191PE
			\pause
			\item 2864 Instâncias -- 16 Atributos
			\pause
			\item Dentre eles: 'text, 'favorited', 'favoriteConunt', 'created', 'ID', 'statusSource', 'screenName', 'retweetCount',
			\pause
			\item 'isRetweet', 'retweeted', dentre outros.		
		\end{itemize}
	\end{exampleblock}
\end{frame}

%++++++++++++++++++++++++++++++++++++++++++++++++++++++++++++++++++++++++++++++++++++++++++++++

%\section{ }
\subsection{ Metodologias de Mineração Dados --- CRISP -- DM}

\begin{frame}\frametitle{ Domínio das técnicas de Mineração de Dados}
	\transdissolve[duration=2, direction=25]
	\begin{figure}[!ht]
		\includegraphics[width=90mm, height=60mm]{Figuras/BigData/RelacaoCrispKddDm.png}
	\end{figure}
\end{frame}


\begin{frame}\frametitle{ Cross Industry Standard Process for Data Maning -- CRISP - DM}
	\transboxout[duration=2, direction=25]
	\begin{figure}[!ht]
		\includegraphics[width=90mm, height=60mm]{Figuras/BigData/CrispDM.png}
	\end{figure}
\end{frame}


\begin{frame}\frametitle{ Etapas CRISP-DM resumo}
	\transdissolve[duration=1, direction=25]
	\begin{figure}[!ht]
		\includegraphics[width=80mm, height=58mm]{Figuras/Crisp/Entendimento.png}
	\end{figure}
\end{frame}

\begin{frame}\frametitle{ Etapas CRISP-DM resumo}
	\transdissolve[duration=1, direction=25]
	\begin{figure}[!ht]
		\includegraphics[width=80mm, height=58mm]{Figuras/Crisp/EntendDados.png}
	\end{figure}
\end{frame}

\begin{frame}\frametitle{ Etapas CRISP-DM resumo}
	\transdissolve[duration=1, direction=25]
	\begin{figure}[!ht]
		\includegraphics[width=80mm, height=58mm]{Figuras/Crisp/PreparaDados.png}
	\end{figure}
\end{frame}

\begin{frame}\frametitle{ Etapas CRISP-DM resumo}
	\transdissolve[duration=1, direction=25]
	\begin{figure}[!ht]
		\includegraphics[width=80mm, height=58mm]{Figuras/Crisp/Model_IA.png}
	\end{figure}
\end{frame}

\begin{frame}\frametitle{ Etapas CRISP-DM resumo}
	\transdissolve[duration=1, direction=25]
	\begin{figure}[!ht]
		\includegraphics[width=80mm, height=58mm]{Figuras/Crisp/Avaliacao.png}
	\end{figure}
\end{frame}

\begin{frame}\frametitle{ Etapas CRISP-DM resumo}
	\transdissolve[duration=1, direction=25]
	\begin{figure}[!ht]
		\includegraphics[width=80mm, height=58mm]{Figuras/Crisp/Implantacao.png}
	\end{figure}
\end{frame}

%\section{KDD e Mineração de Dados}
\subsection{ Minerando os dados do problema}

\begin{frame}\frametitle{ A descoberta de conhecimento (KDD) nas bases da PRF}
	\begin{figure}[!ht]
		%	\centering %\centering % para centralizarmos a figura
		\includegraphics[width=100mm, height=30mm]{Figuras/BigData/Fayyad.png}
		\pause
		\begin{alertblock}{ Diferencça CRISP--DM --- KDD}			
			O CRISP-DM difere do KDD principalmente pelas fases do entendimento do  negócio (anterior ao KDD) e da implantação (posterior ao KDD)
		\end{alertblock}
	\end{figure}
\end{frame}


%++++++++++++++++++++++++++++++++++++++++++++++++++++++++++++++++++++++++++++++++++++++++++++++

%\section{ Tecnologias empregues na Mineração Dados}
\subsection{ Árvore de Decisão}

\begin{frame}\frametitle{ A escolha}
	\begin{exampleblock}{Pontos fortes}
		\begin{itemize}
			\item Custo benefício -- Uma árvore é gerada de maneira simples com resultados abrangente e facilidade de interpretação.
			\pause
			\item Permite fazer predição e classificação intantâneas
		\end{itemize}
	\end{exampleblock}
	\pause
	\begin{alertblock}{Pontos fracos}
		\begin{itemize}
			\item Nó raiz não é facilmente identificável
			\pause
			\item Vários testes até se conseguir resultados satisfatórios		
		\end{itemize}
	\end{alertblock}
\end{frame}


\begin{frame}\frametitle{ Cálculo da Entropia e Entropia Condicional}
	\transdissolve[duration=1, direction=25]
	\centering{Equação da entropia}
	\begin{equation}
	H_{x}=-\sum_{\forall x \in X}P(x)log_{2}P(x)
	\end{equation}
	\centering{Equação da entropia condicional}
	\begin{equation}
	H_{Y|X}= \sum_{x}P(x)H(Y|X=x) =-\sum_{\forall x \in X}P(x) \sum_{\forall y \in Y}P(y|x)log_{2}P(y|x)
	\end{equation}
\end{frame}


\subsection{ Naïve Bayes}

\begin{frame}\frametitle{ Teorema de Bayes e Aprendizagem bayesiana}
	\transdissolve[duration=1, direction=25]
	\centering{Teorema de Bayes: forma tradicional}
	\begin{equation}
	p(C_{k}|x) = \frac{p(_{k})p(x|C_{k})}{p(x)}
	\end{equation}
	
	\centering{Teorema de Bayes: forma simplificada}
	\begin{equation}
	p(posteriori) = \frac{p(priori) * {\text{verossimilhança}}}{\text{evidência}}
	\end{equation}
\end{frame}

\subsection{ Redes Neurais}

\begin{frame}\frametitle{ Classificação com Redes Neurais}
	\transdissolve[duration=1, direction=25]
	Rede Nerual com três camadas identificadas
	\begin{itemize}
		\item Camada de entrada: apresenta-se os padrões à rede
		\item Camada intermediária (ocultas): realiza a maior parte do processamento por conexões ponderadas
		\item Camada de saída: conclusão e apresentação do resultado.
	\end{itemize}
	\begin{figure}[!ht]
		\includegraphics[width=80mm, height=40mm]{Figuras/BigData/camdasIntermediarias.png}
	\end{figure}
\end{frame}

%++++++++++++++++++++++++++++++++++++++++++++++++++++++++++++++++++++++++++++++++++++++++++++++

%\section{ Mineração em Textos}
\subsection{ Tecnologias empregues na Mineração em Textos}

\begin{frame}\frametitle{ TF -- IDF}
	\transdissolve[duration=1, direction=25]
	 \centering{Text Frequence X Inverse Document Frequence: $TF * IDF$}
	\begin{equation}
		idf(term) = \ln(\frac{n_{documents}}{n_{documents\  containing \ term}})
	\end{equation}
	\pause
	\begin{figure}[!ht]
	 		\includegraphics[width=70mm, height=40mm]{Figuras/BigData/ferqPalavrasG.png}
	\end{figure}
\end{frame}

\begin{frame}\frametitle{ Análise de Desempenho aplicados à mineração}

	A área sob a curva AUC (Area under ROC -- Reciver Operating Characeteristic) é uma métrica que determina a qualidade do classificador.
	Ela é calculada cruzando-se verdadeiros positivos com os falso positivos. 
	Esta área varia entre zero e um, quanto mais próximo de 1(um) o classificador conseguiu acertar mais vezes do que errar.
	A seguir a Matriz de Contingência (ou confusão) que agrupa essas atributos.
	%tabela 5
	\begin{table}[ht]
		\centering
		\caption{Matriz de Confusão}
		\vspace{1mm}
		\begin{tabular}{l|c|c}
			\hline
			\textbf{} & \textbf{Predito} & \textbf{}\\
			\hline
			\textbf{Real}  & TP   FN & Positive -- POS\\
			\textbf{Real}  & FP   TN & Negative -- NEG\\
			\hline
			---         & PP   PN &    ---         \\
		\end{tabular}
		\\
		\tiny Fonte: Bradley -- 1997
	\end{table}
\end{frame}

\begin{frame}\frametitle{ Análise de Desempenho aplicados à mineração}
	
	A Matriz Modelo de Confusão sintetiza a Matriz de Confusão
%tabela 6
\begin{table}[ht]
	\centering
	\caption{Matriz modelo de Confusão}
	\vspace{1mm}
	\begin{tabular}{l|c|c}
		\hline
		\textbf{}           & \textbf{Y}     \textbf{$\bar{Y}$}   & \textbf{}\\
		\hline
		\textbf{X}          & P(X,Y)         P(X,$\bar{Y}$)       & Positive -- POS\\
		\textbf{$\bar{X}$}  & P($\bar{X}$,Y) P($\bar{X},\bar{Y}$) & Negative -- NEG\\
		\hline
		---              & P(Y)           P($\bar{Y}$)         &     ---        \\
	\end{tabular}
	\\
	\tiny Fonte: Bradley -- 1997
\end{table}

 Para construir a curva ROC utiliza-se as probabilidades condicionais cruzando-se a taxa de verdadeiros positivos ($ tpr = P(Y|X)$) probabilidade de falsos alarmes ou taxa de falsos positivos será ($ fpr = P(Y,\bar{X})$),
 
 
\end{frame}

\section{Contribuição}


\section{Resultados}

\section{Conclusões}
%++++++++++++++++++++++++++++++++++++++++++++++++++++++++++++++++++++++++++++++++++++++++++++++

\section{Extras}
\subsection{Trabalhos futuros}

\begin{frame}\frametitle{ Aplicativo}
\transboxout[duration=2, direction=25]
	\begin{figure}[ht]
		\centering
		\includegraphics[width=100mm, height=45mm]{Figuras/Extras/InterfaceGrafica.png}
		\caption{ Interface Gráfica Proposta}		
	\end{figure}
\end{frame}

\subsection{Artigos aprovados em eventos} 
\begin{frame}\frametitle{ Artigo aprovado em evento internacional}
	\transboxout[duration=1, direction=25]
	\centering{ANÁLISE DE COMPORTAMENTO RODOVIÁRIO NA REGIÃO METROPOLITANA DO RECIFE-BRASIL -- UMA ABORDAGEM DE MINERAÇÃO DE DADOS}
	\begin{figure}[ht]
		\centering
		\includegraphics[width=80mm, height=25mm]{Figuras/Extras/Cisti122.png}
		\caption{ CISTI -- Conf. Ibérica de Sistemas e Tecnologia da Informação}
	\end{figure}
\end{frame}

\end{document}
